\section{Development}

\subsection{Needs and capacities}

Capacities and needs are synthesized by the same model from independent file pairs: the
capacity effects and thresholds give the capacity distribution, and the need effects and
thresholds give the need distribution. Each resource is modeled as an ordinal outcome with
a cumulative-logit (proportional-odds) model, where $\sigma(x) = 1 / (1 + e^{-x})$ is the
logistic function.

\paragraph{Linear predictor.}
A stratum $s = (v_1, \dots, v_D)$ takes one value from each of the $D$ dimensions of the
effects file, which holds a coefficient $\gamma_{r,d,v}$ for every resource $r$, dimension
$d$ and value $v$. One of the $D$ dimensions is always \texttt{zone\_id}: every stratum
belongs to a zone, and the zone is what agent synthesis later places on the map. Effects
are additive main effects, with no interactions, intercept or reference level:
\begin{equation}
  \beta_{r,s} = \sum_{d=1}^{D} \gamma_{r,d,v_d}.
\end{equation}
If any coefficient is empty, $\beta_{r,s}$ is undefined and the pair $(s, r)$ is omitted:
an empty effect means the stratum does not participate in the resource, which is not the
same as an effect of zero.

\paragraph{Cutpoints.}
Resource $r$ has ordered integer levels $\ell_{r,1} < \dots < \ell_{r,K_r}$ and sorted
cutpoints $\mu_{r,1} \le \dots \le \mu_{r,K_r-1}$ from the thresholds file; the largest
level has none. Only resources present in both files are synthesized.

\paragraph{Distribution.}
The cumulative probability of the outcome being at most level $\ell_{r,k}$ is
\begin{equation}
  F_{r,s}(k) = \sigma\left(\mu_{r,k} - \beta_{r,s}\right),
  \qquad k = 1, \dots, K_r - 1,
\end{equation}
with $F_{r,s}(0) = 0$ and $F_{r,s}(K_r) = 1$, and the probability of each level is the
consecutive difference
\begin{equation}
  p_{r,s}(\ell_{r,k}) = F_{r,s}(k) - F_{r,s}(k-1).
\end{equation}
These are non-negative and sum to one. The distribution is kept whole, one row per stratum,
resource and level. The synthesis is deterministic; agents are sampled from these
distributions in a later step.

\subsection{Supply and demand}

Aggregate supply and demand use the same model, fed by their own effects and thresholds
files, but each distribution is collapsed to its rounded expected value. For a stratum $s$
and resource $r$, with $p_{r,s}$ as in the previous subsection,
\begin{equation}
  A_{r,s} = \operatorname{round}\left(\sum_{k=1}^{K_r} \ell_{r,k} \, p_{r,s}(\ell_{r,k})\right),
\end{equation}
where $A$ is the supply or the demand. The rounding reflects that resources are counted in
whole units. Pairs $(s, r)$ omitted by the model stay empty, which is distinct from a
computed zero, since zero is itself a valid level. The result is wide, one row per stratum
and one column per resource, and deterministic.

\subsection{Agents}

Agents are individuals drawn from each stratum until its supply and demand are used up. A
stratum $s$ is usable only if it has the same dimensions in the supply, demand, capacity and
need files, and a resource $r$ is usable in $s$ only if it is defined in all four. Unusable
strata and resources yield no agents. The zone of $s$ is its \texttt{zone\_id}, with polygon
$\Omega_s$.

\paragraph{Depletion.}
Fix a stratum $s$, and let $a_r^{(0)} = A^{\mathrm{sup}}_{r,s}$ and
$b_r^{(0)} = A^{\mathrm{dem}}_{r,s}$ be its aggregate supply and demand. For agent
$t = 1, 2, \dots$ and each usable resource $r$, the capacity distribution
$p^{\mathrm{cap}}_{r,s}$ is truncated to the levels the remaining supply can still cover
and renormalized,
\begin{equation}
  \tilde p^{\mathrm{cap}}_{r,t}(\ell) =
  \frac{p^{\mathrm{cap}}_{r,s}(\ell) \, \mathbf{1}\left[\ell \le a_r^{(t-1)}\right]}
       {\sum_{\ell' \le a_r^{(t-1)}} p^{\mathrm{cap}}_{r,s}(\ell')},
\end{equation}
and likewise $\tilde p^{\mathrm{need}}_{r,t}$ from $p^{\mathrm{need}}_{r,s}$ truncated at
$b_r^{(t-1)}$. A capacity $x_{r,t} \sim \tilde p^{\mathrm{cap}}_{r,t}$ and a need
$y_{r,t} \sim \tilde p^{\mathrm{need}}_{r,t}$ are drawn independently. If for any $r$ no
level is feasible, the agent cannot be a complete record and generation for $s$ stops.
Otherwise the agent takes $x_{r,t}$ and $y_{r,t}$ for every $r$, and the budgets are
updated,
\begin{equation}
  a_r^{(t)} = a_r^{(t-1)} - x_{r,t}, \qquad b_r^{(t)} = b_r^{(t-1)} - y_{r,t}.
\end{equation}
If an agent draws zero capacity and zero need for every resource, the state is unchanged
and generation stops after it. The number of agents per stratum is the stopping time of
this process rather than an input.

\paragraph{Placement.}
An agent is placed uniformly in $\Omega_s$: a point is drawn uniformly in the polygon's
bounding box and kept if it lies inside $\Omega_s$, and redrawn otherwise.

The result has one row per agent, with its identifier, strata, capacity and need for each
usable resource, and its location, as in the Agents table of the appendix.

\subsection{Desire lines}

Desire lines are individual transactions between pairs of agents: for each resource,
agents with leftover capacity (providers) are matched to agents with leftover need
(consumers) until one side is used up. Resources are handled independently, and only
agents with both a capacity and a need value for a resource take part in it.

\paragraph{Provider.}
Fix a resource $r$, and let $\kappa_i$ and $\nu_i$ be agent $i$'s residual capacity and
need for $r$, initialized from the agents file. While some agent has $\kappa_i > 0$ and
some agent has $\nu_i > 0$, a provider $p$ is drawn from $\{i : \kappa_i > 0\}$ with
probability proportional to $\kappa_p$.

\paragraph{Consumer.}
While $p$ still has capacity left and some agent still has need left, a consumer
$c \neq p$ is drawn from $\{i \neq p : \nu_i > 0\}$ with probability proportional to
\begin{equation}
  \nu_c \, \sigma\!\left(-\frac{d(p,c)}{L}\right),
\end{equation}
where $d(p,c)$ is the distance between the two agents' locations and $L$ is the
bounding-box diagonal of the whole agent cloud, computed once so the decay does not
depend on the coordinate units. Nearer consumers are more likely; the decay is a spatial
weight here, not a probability.

\paragraph{Transaction.}
The traded quantity is
\begin{equation}
  q = \min(\kappa_p, \nu_c),
\end{equation}
after which both $\kappa_p$ and $\nu_c$ are reduced by $q$, and a line is recorded from
$p$ to $c$ carrying $q$. Once $\kappa_p$ reaches zero a fresh provider is drawn, and
one provider can be matched to several consumers in a row before that happens. A weight
is always a residual, so a used-up agent is never drawn again, and every transaction
strictly reduces the capacity or the need left for $r$.

The result has one row per transaction, with the resource, the quantity, the provider's
agent identifier, and a two-point line from the provider's location to the consumer's, as
in the Desire lines table of the appendix. The same pair of agents can appear together
more than once, even for the same resource.
